\input{../tex-inputs/latex-leseansicht-vorspann}

\section[Arthur und Olga Schnitzler an Gustav Schwarzkopf, {[}zwischen 7. und 11.? 5. 1906{]}]{L04398 Arthur und Olga Schnitzler an Gustav Schwarzkopf, {[}zwischen 7. und 11.? 5. 1906{]}}
\nopagebreak\mylabel{L04398v}
\rehead{ }\normalsize\beginnumbering\briefempfaengerindex{Schwarzkopf, Gustav@\textsc{Schwarzkopf, Gustav}!zzzSchnitzler, Olga@\emph{von Olga Schnitzler}!1906-05-111@{{[}zwischen 7. und 11.? 5. 1906{]}}|(be}\briefempfaengerindex{Schwarzkopf, Gustav@\textsc{Schwarzkopf, Gustav}!zzzSchnitzler, Arthur@\emph{von Arthur Schnitzler}!1906-05-111@{{[}zwischen 7. und 11.? 5. 1906{]}}|(be}
\toendnotes[C]{\smallbreak\pagebreak[2]}
\correspDesc{Versand  durch Arthur Schnitzler, Olga Schnitzler im Zeitraum [zwischen 7. und 11.? 5. 1906] in Hinterbrühl
\newline{}Übermittlung  im Zeitraum [zwischen 7. und 12.?] 5. 1906 in Hinterbrühl
\newline{}Erhalt  durch Gustav Schwarzkopf im Zeitraum [zwischen 8. und 12. 5. 1906?] in Tiefer Graben 23}\toendnotes[C]{\smallbreak}
\Standort{DLA, A:Schnitzler, HS.1985.1.1897.}
\physDesc{Bildpostkarte, 276 Zeichen
\newline{}Handschrift Arthur Schnitzler: Bleistift, deutsche Kurrent
\newline{}Handschrift Olga Schnitzler: Bleistift, lateinische Kurrent
\newline{}Versand: Stempel: »\nobreak{}\oindex{Hinterbrühl@\textbf{Hinterbrühl}|pwk}\textcolor{gray}{Hin}terbrü{[}hl{]}, \textcolor{gray}{×}. V. \textcolor{gray}{0}{[}6{]}, VI\nobreak{}«.  }\toendnotes[C]{\smallbreak}\pstart{}{\pb}Herrn \textsc{Gustav
                     Schwarzkopf}\pend{}\pstart{}\textsc{Wien I}\oindex{I., Innere Stadt@\textbf{I., Innere Stadt}|pw}\pend{}\pstart{}\textsc{Tiefer Graben
                  23}\oindex{Wien@\textbf{Wien}!I., Innere Stadt@\textbf{I., Innere Stadt}!Tiefer Graben 23@\textbf{Tiefer Graben 23}|pw}.\pend{}{\bigskip}
\pstart
           \centering{}{\pb}\textcolor{gray}{\textbf{285. – Hinterbrühl. Wagnerstraße\oindex{Wagnerstraße [Hinterbrühl]@\textbf{Wagnerstraße [Hinterbrühl]}|pw}.}}\pend
           \vspace{1em}
\pstart
           \noindent{}{\pb}lieber Guſtav, es iſt fabelhaft \label{K_L04371-1v}\edtext{hübſch hier\oindex{Hinterbrühl@\textbf{Hinterbrühl}|pwv}}{\lemma{\textnormal{\emph{hübsch hier}}}\Cendnote{\textnormal{Zwischen 
                  7. 5. 1906 und 14. 5. 1906 wohnte Schnitzler mit seiner Familie\pwindex{Schnitzler, Olga 17.\,1.\,1882 Wien – 13.\,1.\,1970 Lugano@\textsc{Schnitzler, Olga} (17.\,1.\,1882 Wien – 13.\,1.\,1970 Lugano)|pwkv}\pwindex{Schnitzler, Heinrich 9.\,8.\,1902 Hinterbrühl – 12.\,7.\,1982 Wien@\textsc{Schnitzler, Heinrich} (9.\,8.\,1902 Hinterbrühl – 12.\,7.\,1982 Wien)|pwkv}
                  im Hotel Radetzky\oindex{Hotel Radetzky@\textbf{Hotel Radetzky}|pwk} in Hinterbrühl\oindex{Hinterbrühl@\textbf{Hinterbrühl}|pwk}. Die Abreise erfolgte
                  demnach an einem Montag und nicht, wie auf der vorliegenden Karte noch angedacht, bereits am Sonntag. Das und
                  die Aufforderung an Schwarzkopf\pwindex{Schwarzkopf, Gustav 7.\,11.\,1853 Wien – 13.\,11.\,1939 ebd.@\textsc{Schwarzkopf, Gustav} (7.\,11.\,1853 Wien – 13.\,11.\,1939 ebd.)|pwk}, doch ebenfalls herzukommen, 
                  deutet darauf hin, dass die Karte am Anfang des Aufenthalts verfasst wurde.}}}\label{K_L04371-1}, u das Hotel\oindex{Hotel Radetzky@\textbf{Hotel Radetzky}|pwv} wirklich  gut. Vielleicht ko{\geminationm}en wir nach dem Einſ. Weg\pwindex{Schnitzler, Arthur 15. 5. 1862 Wien – 21. 10. 1931 ebd.@\textsc{Schnitzler, Arthur} (15. 5. 1862 Wien – 21. 10. 1931 ebd.)!einsame Weg. Schauspiel in fünf Akten@\strich\emph{Der einsame Weg. Schauspiel in fünf Akten}|pw}\eventindex{Theater an der Wien@\textbf{Theater an der Wien}!Aufführung von Der einsame Weg, 15.5.1906@Aufführung von Der einsame Weg, 15.5.1906|pwv} wieder heraus. Dem \label{K_L04371-2v}\edtext{Buben\pwindex{Schnitzler, Heinrich 9.\,8.\,1902 Hinterbrühl – 12.\,7.\,1982 Wien@\textsc{Schnitzler, Heinrich} (9.\,8.\,1902 Hinterbrühl – 12.\,7.\,1982 Wien)|pwv} gehts{ }ſchon glänzend}{\lemma{\textnormal{\emph{Buben … glänzend}}}\Cendnote{\textnormal{Heinrich\pwindex{Schnitzler, Heinrich 9.\,8.\,1902 Hinterbrühl – 12.\,7.\,1982 Wien@\textsc{Schnitzler, Heinrich} (9.\,8.\,1902 Hinterbrühl – 12.\,7.\,1982 Wien)|pwk} hatte eine Drüsenschwellung,
                  mit der er seit 1. 5. 1906 in Behandlung war.}}}\label{K_L04371-2}. Wa{\geminationn} er{\pb}ſcheinen Sie?\pend
           
\pstart
           Herzlichſt Ihr{\\[\baselineskip]}\spacefill\mbox{A.}\pend
           \leftskip=0em{}
\pstart
           \noindent{}(So{\geminationn}tag fahr ich \textcolor{gray}{hinein})\pend
           \selectlanguage{ngerman}\vspace{1em}\pstart {[}hs. O. Schnitzler:{]} Die schönſten Grüsse! \spacefill\mbox{OS.}\pend{}\selectlanguage{ngerman}\endnumbering\briefempfaengerindex{Schwarzkopf, Gustav@\textsc{Schwarzkopf, Gustav}!zzzSchnitzler, Olga@\emph{von Olga Schnitzler}!1906-05-071@{{[}zwischen 7. und 11.? 5. 1906{]}}|)be}\briefempfaengerindex{Schwarzkopf, Gustav@\textsc{Schwarzkopf, Gustav}!zzzSchnitzler, Arthur@\emph{von Arthur Schnitzler}!1906-05-071@{{[}zwischen 7. und 11.? 5. 1906{]}}|)be}\mylabel{L04398h}
\begin{anhang}
\end{anhang}\newcommand{\dateiname}{L04398}\newcommand{\titel}{Arthur und Olga Schnitzler an Gustav Schwarzkopf, [zwischen 7. und 11.? 5. 1906]}\newcommand{\editorInnen}{Herausgegeben von Jahnke, SelmaMüller, Martin Anton}\input{../tex-inputs/latex-leseansicht-abspann}
